\section{Performance Evaluation and Cross-Validation}
\label{sec:evaluation}

\textbf{Three metrics, each guarding a different failure.} Every model is scored
on three complementary regression metrics rather than one. $R^2$ reports the
share of outcome variance explained, which is scale-free and so comparable across
the four targets. RMSE gives the typical error in the target's own units and
penalises large misses heavily. MAE gives the plain average absolute error, which
is robust to a few extreme cases. Reporting all three means no single number can
flatter a model: a low RMSE paired with a poor $R^2$, for instance, would expose
a model that looks precise only because the outcome barely varies. Because the
task is regression, class imbalance does not apply, but the target distributions
were still inspected. Both participation outcomes are right-skewed, advanced
Math especially, and that skew reappears in the residuals below.

\textbf{Errors in interpretable units.} For the selected model the grade targets
carry an MAE of 4.11 points (Math) and 3.45 (English) on the 0--100 scale, with
RMSE 5.37 and 4.53. The participation targets carry an MAE of 5.8 and 13.1
percentage points, with RMSE 8.5 and 17.8. The English participation error is the
largest in absolute terms, consistent with that outcome having the widest spread
across schools.

\textbf{Residual diagnostics.} Figure~\ref{fig:residuals} shows the distribution
of the nested out-of-fold residual (actual minus predicted) per target. Because
these predictions come from outer folds never used for selection or tuning, the
diagnostic matches the protocol behind Table~\ref{tab:leaderboard}. Every mean
residual is within $0.03$ of zero, so there is no material \emph{global} bias,
though this does not rule out subgroup bias or calibration error at the extremes.
The grade models are near-symmetric. The participation models are centred on zero
but mildly right-skewed, because the outcome is itself bounded and right-skewed:
the model slightly over-predicts the many low-participation schools and
under-predicts the rare high ones. No panel shows a second peak or heavy
structured tail, though residual-versus-fitted and subgroup diagnostics would be
needed before attributing the remainder to irreducible noise.

\begin{figure}[H]
\centering
\includegraphics[width=0.78\linewidth]{section6/residuals_nested.png}
\caption{Nested out-of-fold residual distributions for the selected model, with
each panel's mean and standard deviation annotated. The coral mean line sits on
the grey zero line in all four panels.}
\label{fig:residuals}
\end{figure}
